% titel: Blatt 10, A1: Seilwelle
% skript: Kap. Wellen, Transversalwelle auf gespanntem Seil

% Seilwelle: Geschwindigkeit c = sqrt(F/mu), Wellenpaket aus Seilenden,
% Ueberlagerung gegenlaeufiger Störungen.

Eine Transversalwelle auf einem gespannten Seil ist das einfachste Beispiel
einer Welle, weil sich ihre Ausbreitungsgeschwindigkeit fast ohne Rechnung
aus einer einzigen Kraft-überlegung ergibt. Um die Aufgabe zu verstehen,
muss man begreifen, woher die Formel $c=\sqrt{F/\mu}$ kommt --- sie ist keine
Definition, sondern das Ergebnis einer dynamischen Überlegung.

\begin{werkzeugkasten}[Warum $c=\sqrt{F/\mu}$? Laengsspannung und Traegheit]
Was bestimmt, wie schnell eine Störung entlang eines Seils läuft? Zwei
gegenläufige Einflüsse:

\begin{itemize}
  \item \textbf{Die Zugkraft $F$ (Spannung):} sie zieht jedes Seilelement
        zurück in seine Ruhelage. Je stärker gespannt das Seil, desto
        schneller ``erholt'' sich eine Auslenkung --- die Welle läuft
        schneller.
  \item \textbf{Die Massenträgheit $\mu=\rho A$ (Masse pro Länge):} ein
        schweres, träges Seil reagiert träge auf die Rückstellkraft --- die
        Welle läuft langsamer.
\end{itemize}

Die Physik steckt in der Kombination: beide hängen gegenläufig, also
\begin{equation}
  c=\sqrt{\frac{F}{\mu}}=\sqrt{\frac{F}{\rho A}}\;.
\end{equation}
Das ist formal dieselbe Form wie $v=\sqrt{D/m}$ des harmonischen Oszillators:
Rückstell-"Größe" ($F$ bzw. Federhärte $D$) über Trägheit ($\mu$ bzw.
Masse $m$). Auch die Schallgeschwindigkeit $c=\sqrt{\kappa p/\rho}$ hat
genau diese Form: eine Art ``Steifigkeit'' über die Dichte.

\textbf{\"Ober ``Faktor Zwei``:} Bei einer Auslenkung wirkt an einem
Seilelement die Seilspannung $F$ an \emph{beiden} Enden; die Rückstellkräfte
addieren sich. Das ist der Ursprung des als Wurzel geschriebenen
Zusammenhangs, nicht eine freie Konstante.

\textbf{Größenordnung:} Gardinen- oder Paketschnur mit $F\approx10$~N,
$\mu\approx10^{-3}$~kg/m ergibt $c\approx100$~m/s --- viel schneller, als man
händisch ``schwingt''. Eine straffe Gitarrensaite (kleines $\mu$, gro"se $F$)
erreicht hunderte m/s.
\end{werkzeugkasten}

\begin{werkzeugkasten}[Ein Wellenpaket am Seilende erzeugen]
Wird das Seilende bei $x=0$ mit $y(0,t)$ ausgelenkt, so wandert diese
Auslenkung unverändert nach rechts:
\begin{equation}
  y(x,t)=y(0,\,t-x/c)\;.
\end{equation}
Jeder ``Momentanwert'' der Handbewegung braucht die Zeit $L/v=x/c$, um zum Ort
$x$ zu gelangen. Solche Funktionen der Kombination $t-x/c$ hei"sen
\emph{vorwärtslaufende Wellen}; eine Funktion von $t+x/c$ läuft nach links.

Ein \emph{Wellenpaket} entsteht, wenn das Ende nur über ein Zeitintervall
$\Delta t$ bewegt wird. Die Störung hat dann räumliche Länge $c\,\Delta t$:
Die Vorderseite wurde zu Beginn, die Rückseite am Ende der Bewegung
ausgesandt und beide laufen mit gleichem $c$ --- die Länge im Raum bleibt
konstant ($c\cdot\Delta t$).
\end{werkzeugkasten}

\begin{werkzeugkasten}[Erkennen von Wellenlänge, Periode, Phase]
In einer Funktion $y=A\sin(\omega t-kx)$ steckt alles:
\begin{itemize}
  \item Kreisfrequenz $\omega$ (Schwingung an \emph{einem} Ort),
        Frequenz $f=\omega/2\pi$.
  \item Wellenzahl $k$ (Schwingung im \emph{Raum} zu einer Zeit),
        Wellenlänge $\lambda=2\pi/k$.
  \item Phasengeschwindigkeit $c=\omega/k=\lambda f$ (die Störung läuft mit
        $x/t=\omega/k$).
  \item Fortschreiten nach rechts: negatives Vorzeichen $-kx$; nach links:
        $+kx$.
\end{itemize}
Maxima sind dort, wo das Argument einem Vielfachen von $\pi/2$ entspricht ---
der ``Wellenberg'' wandert also mit der Phasengeschwindigkeit.
\end{werkzeugkasten}

\begin{achtung}[Grenzen der Näherung]
\begin{itemize}
  \item \textbf{Kleine Auslenkungen} ($y\ll\lambda$, kleine Steigung):
        nur dann ist die Diagonal-Spannung ~$F$ konstant und $c$ unabhängig
        von der Amplitude. Bei grossen Auslenkungen wird die Tangentialkraft
        gr"o"ser und die Welle \emph{nichtlinear} (Wellenkämme laufen den
        Tälern davon).
  \item \textbf{Ideales, masseloses Spannungsmodell:} Biegesteifigkeit,
        innere Reibung und Schwerkraft (hängendes statt liegendes Seil) sind
        vernachl"assigt. Ein senkrecht hängendes Seil hätte zum Beispiel
        wegen des Gewichts eine ortsabhängige Spannung und damit
        ortsabhängiges $c$.
  \item Die Formel $c=\sqrt{F/\mu}$ gilt nur für die \emph{transversale}
        Ausbreitung entlang eines gespannten Seils, nicht für Schall in
        Gasen oder Körperwellen.
\end{itemize}
\end{achtung}
